% =====================================================================
%  BLOCS DE MISE A JOUR DU RAPPORT — a coller dans le .tex existant
%  (chaque bloc indique OU il remplace / s'insere)
% =====================================================================


% ---------------------------------------------------------------------
% BLOC 1 — remplace le paragraphe final de la sous-section 7.1
%          "Variational Objective" (celui qui se termine par
%          "...Thus, we can use Moreau's formula to get a new expression")
% ---------------------------------------------------------------------

The velocity field is then recovered via
\begin{equation}
    v_\theta(x_t, t) = \frac{\hat{x} - x_t}{1-t}.
\end{equation}
Two numerical subtleties arise from this denoising reduction. First, the
rescaling $z = x_t/t$ of Algorithm~\ref{Prox_for_vector_field} blows up as
$t \to 0$. In practice we therefore warm-start every unrolled network directly
from $z = x_t$ (i.e.\ we fold the $1/t$ factor into the learnable operators
rather than applying it explicitly), which is well defined on all of
$[0,1]$. Second, the final division by $1-t$ blows up as $t \to 1$; at
inference we clamp the denominator, $v_\theta(x_t,t) = (\hat{x}-x_t)/\max(1-t,\varepsilon)$
with $\varepsilon = 5\cdot 10^{-4}$. At training time the same $t\to 1$
ill-conditioning is handled analytically by the loss described in
Section~\ref{subsubsec:xpred}.


% ---------------------------------------------------------------------
% BLOC 2 — NOUVELLE sous-section, a inserer dans 7.2 "Learning Strategies"
%          (par ex. juste apres 7.2.3 "Choice of Proximal Operator")
% ---------------------------------------------------------------------

\subsubsection{Output parameterisation: $x$-prediction with a velocity-space loss}
\label{subsubsec:xpred}

A model estimating the vector field can be parameterised in two equivalent ways.
In \emph{velocity prediction} ($v$-pred) the network outputs the velocity
directly, $v_\theta(x_t,t)$, and is trained with the standard CFM loss
$\|v_\theta - (x_1-x_0)\|^2$. In \emph{clean-signal prediction} ($x$-pred) the
network instead outputs an estimate $D_\theta(x_t,t) \approx x_1$ of the clean
image, and the velocity is reconstructed \emph{a posteriori} via
$v_\theta(x_t,t) = (D_\theta(x_t,t) - x_t)/(1-t)$. This is the natural output of
our unrolled solvers, whose primal iterate $\hat{x}$ is by construction an
estimate of $x_1$ (the MAP denoiser of Section~\ref{sec:denoising}); we flag such
models with the attribute \texttt{predicts\_x1}.

The two parameterisations are made \emph{exactly} equivalent by the choice of
loss. Using $x_1 - x_t = (1-t)(x_1 - x_0)$, the velocity error factorises as
\begin{equation}
    \bigl\| v_\theta(x_t,t) - (x_1-x_0) \bigr\|^2
    = \frac{1}{(1-t)^2}\,\bigl\| D_\theta(x_t,t) - x_1 \bigr\|^2 .
    \label{eq:vloss_equiv}
\end{equation}
Hence, rather than regressing $D_\theta$ onto $x_1$ with a plain
$\|D_\theta - x_1\|^2$ objective, we train it with the \emph{velocity-space
weighted} loss
\begin{equation}
    \mathcal{L}_{\text{x-pred}}(\theta)
    = \mathbb{E}_{t, x_0, x_1}\!\left[
        w(t)\,\bigl\| D_\theta(x_t,t) - x_1 \bigr\|^2
    \right],
    \qquad
    w(t) = \frac{1}{\max\!\bigl((1-t)^2,\ \varepsilon^2\bigr)},
    \label{eq:xpred_loss}
\end{equation}
with $\varepsilon = 0.05$. By \eqref{eq:vloss_equiv}, up to the clamp this is
\emph{identical} to the CFM objective $\mathbb{E}\|v_\theta-(x_1-x_0)\|^2$: the
network is optimised for the correct velocity-matching goal, while regressing a
target ($x_1$) that is bounded, well scaled and image-like at all times. The
clamp merely caps the weight in the immediate neighbourhood of $t=1$, where the
raw $1/(1-t)^2$ factor would otherwise diverge, trading a negligible bias for
stability.

This parameterisation is not cosmetic. A plain $x$-prediction loss
($w \equiv 1$) makes the regression target $x_1$ nearly independent of $t$ and,
for a low-variance data distribution (e.g.\ a single MNIST class), produces a
weak, poorly-conditioned gradient: the network minimises the loss by collapsing
onto the \emph{data mean} and ignoring its input $x_t$, yielding a constant
generator. The $1/(1-t)^2$ weighting (equivalently, matching in velocity space)
restores the informative high-$t$ regime and removes this collapse. This is the
$x$-prediction / $v$-loss formulation advocated in the ``back-to-basics'' line of
work on diffusion and flow parameterisations, and it is the setting used for all
MNIST experiments below.


% ---------------------------------------------------------------------
% BLOC 3 — REMPLACE integralement la sous-section 7.4.3
%          "Convolutional UNNs (per-layer weights)"
% ---------------------------------------------------------------------

\subsubsection{Convolutional UNNs (per-layer weights)}

For image data the dense operator $W_k$ is replaced by a learnable 2D
convolution and $V_k$ by a transposed convolution of matching shape. Concretely,
$W_k$ maps the single-channel primal image to a $d_u = \texttt{ic}$-channel dual
feature space with a $9\times 9$ kernel and padding $4$ (so that
$\texttt{conv\_transpose2d}$ is the exact adjoint of $\texttt{conv2d}$ and spatial
size is preserved), where the number of intermediate channels $\texttt{ic}$ is a
hyperparameter. The convolution weights are Kaiming-initialised. Unless stated
otherwise the proximal operator is either the analytical $\ell_1$ dual prox
(\texttt{L1ProxConv}, a per-channel clipping with a learned time-dependent radius,
see Section~\ref{subsubsec:prox_choice}) or a \texttt{DoubleConvTime} block (two
convolutional layers with GroupNorm, SiLU and additive time injection). Three
families are provided:
\begin{itemize}
    \item \texttt{ConvDFB\_UNN}: convolutional DFB, LFO and LNO.
    \item \texttt{ConvDiFB\_UNN}: convolutional DiFB, LFO and LNO.
    \item \texttt{ConvScCP\_UNN}: convolutional ScCP, LFO and LNO.
\end{itemize}


% ---------------------------------------------------------------------
% BLOC 4 — NOUVELLE sous-section, a inserer dans 7.4 "Model Variants"
%          (par ex. apres 7.4.3), OU juste avant les resultats MNIST.
% ---------------------------------------------------------------------

\subsubsection{Breaking the odd symmetry of the proximal velocity field}
\label{subsubsec:oddsym}

The convolutional ScCP block with the analytical $\ell_1$ dual prox has a
structural property that is invisible on the (centred, near-symmetric) two-moons
data but harmful on images: as a function of its input, the induced velocity
field is \emph{odd}. Indeed, with warm start $z = x_t$, dual initialisation
$u_0 = 0$, and no convolutional bias, every ScCP update
(primal step, extrapolation, dual forward step) is \emph{linear} in $(x,u,z)$,
and the symmetric clipping prox satisfies $\mathrm{prox}(-a) = -\mathrm{prox}(a)$.
By induction the map $x_t \mapsto \hat{x}$ is odd, hence
\begin{equation}
    v_\theta(-x_t, t) = -\,v_\theta(x_t, t) .
    \label{eq:oddness}
\end{equation}
This holds \emph{architecturally}, even at initialisation and before any
training; we verify it numerically ($\|v_\theta(-x_t)+v_\theta(x_t)\|/\|v_\theta(x_t)\|
\approx 10^{-6}$).

Property \eqref{eq:oddness} propagates through the generative ODE:
$\mathrm{gen}(-x_0) = -\mathrm{gen}(x_0)$. On MNIST, a digit has a black
background ($\approx -1$) and white strokes, so its negation is a
\emph{colour-inverted} image (white background). Since the noise $x_0$ and $-x_0$
are equally likely, the generator is forced to produce a distribution that is
symmetric about $0$: roughly half of the generated samples come out
colour-inverted, and the model cannot represent the asymmetric MNIST
distribution. This is a genuine limitation inherited from the odd symmetry of the
underlying convex/proximal problem.

We break the symmetry with two lightweight, \emph{zero-initialised} mechanisms,
so that training always \emph{starts} from the known-working symmetric model and
only departs from it if beneficial:
\begin{enumerate}
    \item \textbf{Learned convolutional bias (\texttt{w\_bias}).} A per-dual-channel
          bias $b$ is added in the dual forward operator,
          $\sigma_k\,(W_k y_{k+1} + b)$, which amounts to penalising
          $g(W_k\,\cdot - b)$ instead of $g(W_k\,\cdot)$. The shift makes the map
          non-odd in $x_t$. With $b$ zero-initialised the model starts exactly at
          the symmetric operator.
    \item \textbf{Asymmetric dual prox.} The symmetric clip
          $\mathrm{clip}(u,-r,r)$ is replaced by a two-threshold clip
          $\mathrm{clip}(u,-\lambda^-,\lambda^+)$ with independent
          $\lambda^\pm = \mathrm{softplus}(\text{base} \pm \text{delta}) \ge 0$.
          The interval still contains $0$, so this remains the exact proximal
          operator of the (asymmetric but convex) penalty
          $\lambda^+ \max(x,0) + \lambda^- \max(-x,0)$; it is non-odd as soon as
          $\lambda^+ \neq \lambda^-$. With \text{delta} zero-initialised,
          $\lambda^+ = \lambda^-$ and the model again starts symmetric.
\end{enumerate}
Both fixes preserve the proximal/variational interpretation of the layer (they
correspond to a shifted argument and to an asymmetric convex penalty
respectively) while removing the parity constraint that made unmodified ScCP
unable to fit image data.


% ---------------------------------------------------------------------
% BLOC 5 — REMPLACE la sous-section 7.8 "Results on MNIST"
%          (intro + "Experimental setting"), la Table 2 et les figures
%          restant inchangees. Ajoute la remarque "ce qui n'a pas marche".
% ---------------------------------------------------------------------

\subsection{Results on MNIST}
\label{subsec:mnist}

\paragraph{From two moons to images.} Moving to the $28\times 28$ image space of
MNIST requires adapting the architectural inductive biases: dense operators scale
poorly and ignore spatial locality. As described in
Section~\ref{subsubsec:convolutional}, we replace the linear operator $W_k$ by a
convolution mapping the $1$-channel primal image to an $\texttt{ic}$-channel dual
feature space ($9\times 9$ kernel, padding $4$). Among the many variants of
Table~\ref{tab:variants}, our MNIST study focuses on a \textbf{single} model
family, the one that proved both the most faithful to the variational picture and
the most competitive: the convolutional strongly-convex Chambolle--Pock network
\texttt{ConvScCP\_UNN} with the \textbf{analytical $\ell_1$ dual prox}, trained in
the \textbf{$x$-prediction / velocity-loss} setting of
Section~\ref{subsubsec:xpred} and the \textbf{LNO} regime, operating directly in
pixel space. The proximal operator is thus fully interpretable (a per-channel
$\ell_1$ clipping with learned time-dependent radius), and the only black-box
component is the linear analysis operator $W_k$.

\paragraph{Experimental setting.} We train \texttt{ConvScCP\_UNN} on the
single-class subset of MNIST restricted to the digit $0$ (independent coupling,
batch size $128$, $50$ epochs, Adam). We sweep the number of unrolled iterations
$K \in \{5,10,20\}$ and the number of intermediate channels
$\texttt{ic} \in \{32,64,128\}$, and compare against \texttt{SmallUNet}
baselines of comparable or larger parameter count trained under the same Flow
Matching objective. Beyond the training loss we report two generation-quality
proxies computed on generated samples: \texttt{std\_gen}, the empirical
per-pixel standard deviation across the generated batch (a diversity proxy), and
the anisotropic gradient energy
\[
    \mathrm{GE}(x) = \frac1N\sum_{i,j}\lvert x_{i+1,j}-x_{i,j}\rvert
                   + \frac1N\sum_{i,j}\lvert x_{i,j+1}-x_{i,j}\rvert ,
\]
the mean $\ell_1$ total variation of the batch (a sharpness / high-frequency
proxy), both compared to their reference value on true MNIST zeros,
$\mathrm{GE}_{\text{reel}} = 0.2998$. Both statistics are computed on a small
batch ($n_{\text{cell}} = 4$) and should be read as indicative rather than
conclusive.

% (Table 2 et Figure 4 inchangees)

\paragraph{Effect of $K$ and \texttt{ic}.} Increasing $K$ from $5$ to $10$
generally improves both the final loss and the proximity of GE to
$\mathrm{GE}_{\text{reel}}$: at $\texttt{ic}=128$, GE rises from $0.291$ ($K=5$)
to $0.329$ ($K=10$) and $0.335$ ($K=20$), essentially matching the real-data
value. Visually this corresponds to the disappearance of the high-frequency
speckle seen at $K=5$ and to increasingly closed, well-formed contours. For fixed
$K$, increasing \texttt{ic} from $32$ to $128$ consistently improves both loss
and GE, confirming that the richer dual feature space helps the shared proximal
block resolve finer detail. The best configuration, $K=20$, $\texttt{ic}=128$
($210\,440$ parameters), attains a gradient-energy statistic closer to
$\mathrm{GE}_{\text{reel}}$ than \emph{any} \texttt{SmallUNet} baseline, including
the $2.8$M-parameter one. The gain from $K=10$ to $K=20$ is however
non-monotonic in the training loss (e.g.\ at $\texttt{ic}=32$ it rises from
$0.280$ to $0.332$), consistent with the known sensitivity of deep unrolled
schemes to vanishing/exploding signals across iterations.

\paragraph{What did not work.} Several negative findings shaped the final setup
above and are worth recording.
\begin{itemize}
    \item \textbf{Optimal-transport coupling breaks ScCP, but not the UNet
          baseline.} Minibatch OT couplings \cite{tong2023improving}
          (\texttt{ExactOptimalTransportConditionalFlowMatcher}) straighten the
          paths and help on the $2$-D two-moons problem. On MNIST they train the
          \texttt{SmallUNet} baseline without difficulty, yet they \emph{destroy}
          the \texttt{ConvScCP\_UNN} field, which degenerates to a blurry,
          near-constant velocity. The failure is therefore \emph{not} a generic
          high-dimensional OT collapse but an interaction with the constrained
          proximal architecture: unlike the free-form UNet, the odd, low-capacity
          ScCP operator (Section~\ref{subsubsec:oddsym}) cannot represent the
          strongly input-dependent, asymmetric noise-to-data routing that OT
          coupling imposes on each minibatch, and averages it out. All ScCP MNIST
          experiments therefore use independent coupling, while OT remains usable
          for the UNet baseline.
    \item \textbf{Odd symmetry / colour inversion.} As analysed in
          Section~\ref{subsubsec:oddsym}, the plain $\ell_1$ ConvScCP is an odd
          operator, so about half the generated samples came out colour-inverted
          and the generated distribution was symmetric about $0$. This is
          architectural, not an optimisation artefact; it is removed only by the
          \texttt{w\_bias} or asymmetric-prox mechanisms.
    \item \textbf{$x$-prediction collapse without velocity weighting.} A plain
          $\|D_\theta - x_1\|^2$ loss led the network to collapse onto the data
          mean (a constant generator), especially on a single low-variance class.
          The velocity-space weighting \eqref{eq:xpred_loss} was necessary to
          recover a data-dependent field.
    \item \textbf{Latent-space Flow Matching blurred the samples and stalled
          ScCP.} We also tried running the flow in the latent space of a frozen,
          lightly-variational convolutional autoencoder: each $28\times 28$ image
          is encoded to a $(c_{\text{lat}},7,7)$ code, the flow is matched
          directly on the codes (\texttt{LatentScCP\_UNN}), and the AE decodes
          only for visualisation. This consistently produced \emph{blurrier}
          digits than the pixel-space model and, more tellingly, the ScCP
          iterations \emph{converged poorly}. The likely cause is that the
          analytical $\ell_1$ dual prox is no longer an appropriate prior in the
          latent domain: sparsity is a sensible regulariser on a
          (convolutional/wavelet) pixel representation, but the AE latent codes
          are dense, entangled features for which $\ell_1$ clipping is not a
          meaningful proximal step, so the unrolled solver loses the variational
          structure it relies on. All reported experiments are therefore in pixel
          space.
    \item \textbf{ScCP does not beat DFB despite stronger guarantees}, and
          \textbf{shared (tied) weights underperform per-layer weights}, as
          already observed on two moons (Section~\ref{subsec:2moons}); neither
          gap closes on MNIST.
\end{itemize}

\paragraph{Next steps.} (1) Compute a proper generation-quality metric (FID or a
Wasserstein distance) to disambiguate diversity from sharpness. (2) Add training-loss
curves \texttt{ConvScCP\_UNN} vs.\ \texttt{SmallUNet} across epochs, to tell
whether the non-monotonic loss at large $K$ stems from slower convergence or
instability. (3) Extend the digit-$0$ study to the multi-class setting (all ten
digits) to check whether the $x$-prediction / $v$-loss fix removes the collapse
globally or only between classes.
